% !PRINTED 568
\weberpage{568}
%<<CONT
die Summe von zwei Quadraten reeller Linearformen. Das Product
 $\eta \eta'$ tritt dann mit dem Factor 2 in der Function $f$ auf.
Daher ist $f$ unter allen Umst\"anden eine positive Form.

Nach (7), (8) l\"asst sich aber $f$ auch als quadratische Form
der Variablen $\eta_1, \eta_2, \ldots, \eta_n$
\weq{9}{f = F(\eta_1, \eta_2, \ldots, \eta_n) = \sum^{i,k} a_{i,k}\,\eta_i \eta_k}
auffassen, und hierin ist, wenn $\eta_i$ reell ist, $a_{i,i} = 1$, und wenn
$\eta_i, \eta_k$ ein conjugirt imagin\"ares Paar ist, $a_{i,k} = 1$. Alle \"ubrigen
Co\"efficienten sind $= 0$. Die Determinante von $F$ ist also $= \pm 1$,
weil in jeder Zeile und in jeder Colonne nur ein von Null verschiedenes
 Element mit dem Werthe 1 vorkommt. Nach Bd. I,
\S.\,56 ist die Determinante der Form $f$ gleich dem Product der
Determinante von $F$ und dem Quadrate der Determinante der
linearen Formen $\eta$, also nach (5)
\weq{10}{= \pm [N_a(\varphi)]^2 \Delta,}
wo das Zeichen so zu w\"ahlen ist, dass der Werth positiv ausf\"allt.

Nach \S.\,155, (24) kann man nun f\"ur die $x_i$ solche ganze
rationale Zahlen setzen, dass
\weq{11}{f(x_1, x_2, \ldots, x_n) < \Theta_n\, n \sqrt[n]{\pm \Delta N_a(\varphi)^2}}
wird, worin $\Theta_n$, wenn $n$ gr\"osser als 1 ist, ein echter Bruch und
immer kleiner als $\sqrt[n]{0{,}404}$ ist. Setzt man diese Werthe von $x_i$ in
(4) ein, so geht $\eta$ in eine ganze durch $\varphi$ theilbare Zahl $\beta$ \"uber,
und wir haben also damit den Satz:

2. \gesperrt{Ist $\varphi$ irgend ein ganzes Functional in $\Omega$, so giebt}
\gesperrt{es eine durch $\varphi$ theilbare ganze Zahl $\beta$ in $\Omega$ von}
\gesperrt{der Eigenschaft, dass}
\weq{12}{S\,|\eta|^2 < \Theta_n\, n \sqrt[n]{\pm \Delta [N_a(\varphi)]^2}}
\gesperrt{wird, worin $\Delta$ die Grundzahl des K\"orpers und}
\gesperrt{$\pm \Delta$ positiv ist.}

Das Product der conjugirten Werthe $\eta$ ist gleich der Norm
von $\eta$, und daher ist nach dem H\"ulfssatze 1. [mit R\"ucksicht auf
(7) und (8)]
\weq{13}{n \sqrt[n]{[N(\eta)]^2} \leqq S\,|\eta|^2,}
und also folgt nach (12) die auf die absoluten Werthe bez\"ugliche
 Ungleichung:
\weq{14}{N(\eta)^2 < \Theta_n^n\, [N_a(\varphi)]^2\, \Delta.}
Da nun $\eta$ durch $\varphi$ theilbar ist, so k\"onnen wir
\[ \varepsilon \eta = \varphi \psi, \quad \pm N(\eta) = N_a(\varphi)\, N_a(\psi) \]
%CONT>>
